\section{PENDAHULUAN}

\subsection{Konteks Masalah}

Indonesia merupakan salah satu negara penghasil sampah makanan terbesar di dunia. Berdasarkan data dari Badan Pusat Statistik (BPS) tahun 2024 \cite{bps2024lingkungan}, Indonesia menghasilkan sekitar 23-48 juta ton sampah makanan per tahun \cite{indonesiapolicy2025}, yang setara dengan 115-237 kg per kapita. Angka ini menjadikan Indonesia peringkat kedua di dunia setelah Arab Saudi dalam hal pemborosan makanan per kapita.

Sektor UMKM kuliner menjadi salah satu kontributor terbesar terhadap \textit{food waste}. Data Kementerian Koperasi dan Usaha Kecil dan Menengah \cite{kemenkopukm2024} menunjukkan bahwa UMKM kuliner di Indonesia mencapai lebih dari 5 juta unit usaha, dengan tingkat kerugian akibat makanan surplus yang tidak terjual berkisar antara 10-30\% dari total produksi harian \cite{donatfitri2026}. Kerugian ini tidak hanya berdampak pada aspek finansial UMKM, tetapi juga menciptakan masalah lingkungan yang serius.

\begin{table}[htbp]
\centering
\caption{Data Sampah Makanan di Indonesia}
\label{tab:data_sampah}
\begin{tabularx}{\textwidth}{|l|X|}
\hline
\textbf{Indikator} & \textbf{Data} \\
\hline
Total sampah makanan per tahun & 23-48 juta ton \\
\hline
Sampah per kapita per tahun & 115-237 kg \\
\hline
Kontribusi UMKM kuliner & 10-30\% dari produksi \\
\hline
Persentase sampah organik & 50-70\% total sampah \\
\hline
Emisi CO$_2$ dari food waste & $\sim$50 juta ton per tahun \\
\hline
\end{tabularx}
\end{table}

\subsection{Akar Masalah}

Berdasarkan analisis yang telah dilakukan, terdapat beberapa akar masalah utama yang menyebabkan tingginya \textit{food waste} di sektor UMKM kuliner:

\begin{enumerate}[leftmargin=*]
    \item \textbf{Minimnya Kesadaran Lingkungan} : Banyak pelaku UMKM yang belum menyadari dampak lingkungan dari pembuangan makanan berlebih, sehingga kurang termotivasi untuk mengurangi \textit{food waste}.
    
    \item \textbf{Ketidakpastian Permintaan} : UMKM kuliner sering kali kesulitan memprediksi jumlah permintaan konsumen dari hari ke hari, sehingga cenderung memproduksi makanan dalam jumlah berlebih untuk mengantisipasi lonjakan permintaan.
    
    \item \textbf{Keterbatasan Infrastruktur Penyimpanan} : Banyak UMKM kuliner yang tidak memiliki fasilitas penyimpanan yang memadai untuk menjaga kesegaran makanan, sehingga makanan yang tidak terjual harus segera dibuang.
    
    \item \textbf{Kurangnya Akses ke Pasar Alternatif} : Tidak adanya platform yang menghubungkan UMKM dengan konsumen potensial untuk menjual makanan surplus dengan harga diskon menjadi hambatan utama.
\end{enumerate}

\subsection{Dampak jika Tidak Diselesaikan}

Jika masalah \textit{food waste} di sektor UMKM kuliner tidak segera ditangani, konsekuensi yang akan terjadi dalam 1-2 tahun ke depan meliputi:

\begin{itemize}[leftmargin=*]
    \item \textbf{Kerugian Ekonomi yang Meningkat} : UMKM akan terus kehilangan pendapatan akibat makanan yang terbuang, yang diperkirakan mencapai Rp 500 ribu hingga Rp 2 juta per hari untuk satu unit UMKM kuliner.
    
    \item \textbf{Peningkatan Emisi Gas Rumah Kaca} : Sampah makanan yang terurai di Tempat Pembuangan Akhir (TPA) menghasilkan metana, gas rumah kaca yang 25 kali lebih potensial dari CO$_2$ dalam mempercepat perubahan iklim.
    
    \item \textbf{Beban Lingkungan yang Semakin Berat} : Volume sampah organik yang terus meningkat akan mempercepat kerusakan TPA dan meningkatkan risiko pencemaran tanah serta air tanah.
\end{itemize}

\subsection{Mengapa Sekarang?}

Beberapa faktor yang membuat penyelesaian masalah ini menjadi sangat mendesak saat ini:

\begin{enumerate}[leftmargin=*]
    \item \textbf{Kebijakan Pemerintah} : Pemerintah Indonesia telah menetapkan target pengurangan sampah makanan sebesar 50\% pada tahun 2030 \cite{unsdg2024}, sejalan dengan Sustainable Development Goals (SDG) 12.3 \cite{fao2024state}.
    
    \item \textbf{Kesadaran Konsumen yang Meningkat} : Generasi milenial dan Gen Z menunjukkan peningkatan kesadaran terhadap isu lingkungan, termasuk \textit{food waste}, sehingga membuka peluang pasar untuk solusi yang ramah lingkungan.
    
    \item \textbf{Perkembangan Teknologi} : Kemajuan teknologi \textit{web} dan \textit{mobile application} memungkinkan pengembangan platform yang mudah diakses oleh UMKM dengan biaya yang terjangkau.
\end{enumerate}
